\section{Conceptos fundamentales de estadística inferencial}

En esta sección presentaremos los conceptos básicos que son esenciales para comprender la estadística inferencial. Estos conceptos nos proporcionarán el lenguaje y las herramientas necesarias para realizar inferencias sobre poblaciones a partir de muestras.

\subsection{Estimadores y sus propiedades}

Un \emph{estimador} es un estadístico que se utiliza para estimar el valor de un parámetro poblacional desconocido. Denotamos un estimador del parámetro $\theta$ como $\hat{\theta}$.

\begin{definicion}[Estimador insesgado]
	Un estimador $\hat{\theta}$ es \emph{insesgado} para el parámetro $\theta$ si
	\begin{align}
		E(\hat{\theta}) = \theta
	\end{align}
	
	Es decir, el valor esperado del estimador es igual al verdadero valor del parámetro.
\end{definicion}

\begin{ejemplo}
	La media muestral $\bar{X}$ es un estimador insesgado de la media poblacional $\mu$:
	\begin{align}
		E(\bar{X}) = E\left(\frac{1}{n}\sum_{i=1}^n X_i\right) = \frac{1}{n}\sum_{i=1}^n E(X_i) = \frac{1}{n} \cdot n\mu = \mu
	\end{align}
\end{ejemplo}

\begin{observacion}
	La varianza muestral $S^2 = \frac{1}{n-1}\sum_{i=1}^n (X_i - \bar{X})^2$ es un estimador insesgado de la varianza poblacional $\sigma^2$. El uso de $n-1$ en lugar de $n$ (conocido como corrección de Bessel) es precisamente lo que hace que el estimador sea insesgado.
\end{observacion}

\subsection{Distribuciones muestrales}

Cuando tomamos múltiples muestras del mismo tamaño de una población y calculamos un estadístico para cada muestra, la distribución de estos estadísticos se denomina \emph{distribución muestral}.

\begin{definicion}[Distribución muestral]
	La \emph{distribución muestral} de un estadístico es la distribución de probabilidad de ese estadístico cuando se calcula sobre todas las posibles muestras de tamaño $n$ de una población dada.
\end{definicion}

\begin{ejemplo}
	Si $X_1, X_2, \ldots, X_n$ son variables aleatorias independientes e idénticamente distribuidas con media $\mu$ y varianza $\sigma^2$, entonces la media muestral $\bar{X}$ tiene:
	\begin{align}
		E(\bar{X}) &= \mu \\
		\text{Var}(\bar{X}) &= \frac{\sigma^2}{n}
	\end{align}
	
	Esto significa que la variabilidad de la media muestral disminuye a medida que aumenta el tamaño de la muestra.
\end{ejemplo}

\subsection{Error estándar}

El \emph{error estándar} de un estadístico es la desviación estándar de su distribución muestral. Mide la precisión de la estimación.

\begin{definicion}[Error estándar de la media]
	El error estándar de la media muestral es:
	\begin{align}
		SE(\bar{X}) = \frac{\sigma}{\sqrt{n}}
	\end{align}
	
	Cuando $\sigma$ es desconocida, se estima mediante:
	\begin{align}
		SE(\bar{X}) \approx \frac{s}{\sqrt{n}}
	\end{align}
	donde $s$ es la desviación estándar muestral.
\end{definicion}

\subsection{Ley de los grandes números}

La \emph{ley de los grandes números} establece que, bajo ciertas condiciones, el promedio de los resultados obtenidos de un gran número de ensayos debe estar cerca del valor esperado, y tenderá a acercarse más a medida que se realicen más ensayos.

\begin{teorema}[Ley de los grandes números]
	Si $X_1, X_2, \ldots, X_n$ son variables aleatorias independientes e idénticamente distribuidas con media $\mu$, entonces para cualquier $\epsilon > 0$:
	\begin{align}
		\lim_{n \to \infty} P(|\bar{X}_n - \mu| < \epsilon) = 1
	\end{align}
	
	Es decir, la media muestral converge en probabilidad a la media poblacional.
\end{teorema}

\begin{observacion}
	La ley de los grandes números justifica el uso de promedios muestrales para estimar medias poblacionales. Nos dice que con muestras suficientemente grandes, nuestra estimación será muy cercana al valor verdadero.
\end{observacion}

\subsection{Teorema del límite central}

El \emph{teorema del límite central} (TLC) es uno de los resultados más importantes en estadística. Establece que, bajo ciertas condiciones, la distribución de la suma (o promedio) de un gran número de variables aleatorias independientes se aproxima a una distribución normal, independientemente de la distribución de las variables originales.

\begin{teorema}[Teorema del límite central]
	Si $X_1, X_2, \ldots, X_n$ son variables aleatorias independientes e idénticamente distribuidas con media $\mu$ y varianza $\sigma^2 < \infty$, entonces cuando $n$ es grande:
	\begin{align}
		\frac{\bar{X} - \mu}{\sigma/\sqrt{n}} \xrightarrow{d} N(0,1)
	\end{align}
	
	Es decir, la distribución de la media muestral estandarizada se aproxima a una distribución normal estándar.
\end{teorema}

\begin{observacion}
	En la práctica, el TLC funciona bien cuando $n \geq 30$, aunque esto depende de la forma de la distribución original. Para distribuciones muy asimétricas o con colas pesadas, puede requerirse un $n$ mayor.
\end{observacion}

\begin{ejemplo}
	Supongamos que el tiempo de espera en una fila tiene una distribución exponencial con media $\mu = 5$ minutos. Si tomamos una muestra de $n = 50$ personas, por el TLC, la distribución del tiempo de espera promedio muestral será aproximadamente normal:
	\begin{align}
		\bar{X} \sim N\left(5, \frac{5^2}{50}\right) = N(5, 0.5)
	\end{align}
\end{ejemplo}

\subsection{Importancia del TLC en inferencia estadística}

El teorema del límite central es fundamental para la inferencia estadística porque:

\begin{enumerate}
	\item Nos permite usar la distribución normal para hacer inferencias sobre medias poblacionales, incluso cuando la población original no es normal.
	
	\item Justifica el uso de pruebas Z y t para medias, que estudiaremos en las siguientes secciones.
	
	\item Proporciona la base teórica para la construcción de intervalos de confianza.
	
	\item Explica por qué muchas variables en la naturaleza siguen distribuciones aproximadamente normales (son el resultado de la suma de muchos efectos pequeños e independientes).
\end{enumerate}

\subsection{Resumen de conceptos clave}

\begin{itemize}
	\item \textbf{Estimador:} Estadístico utilizado para estimar un parámetro poblacional.
	\item \textbf{Estimador insesgado:} Estimador cuyo valor esperado es igual al parámetro.
	\item \textbf{Distribución muestral:} Distribución de un estadístico sobre todas las muestras posibles.
	\item \textbf{Error estándar:} Desviación estándar de la distribución muestral.
	\item \textbf{Ley de los grandes números:} La media muestral converge a la media poblacional.
	\item \textbf{Teorema del límite central:} La distribución de la media muestral se aproxima a la normal.
\end{itemize}

Estos conceptos forman la base sobre la cual construiremos las técnicas de estimación y pruebas de hipótesis en las siguientes secciones.
